% ============================================================
% Aegis OS — Final Project Report, AID 325
% Compile with: pdflatex Final_Project_Report_AID325.tex  (twice)
% ============================================================
\documentclass[12pt,a4paper]{article}

\usepackage[T1]{fontenc}
\usepackage[utf8]{inputenc}
\usepackage{geometry}
\usepackage{setspace}
\usepackage{titlesec}
\usepackage{fancyhdr}
\usepackage{hyperref}
\usepackage{booktabs}
\usepackage{longtable}
\usepackage{array}
\usepackage{listings}
\usepackage{xcolor}
\usepackage{colortbl}
\usepackage{graphicx}
\usepackage{caption}
\usepackage{enumitem}
\usepackage{amsmath}
\usepackage{microtype}
\usepackage{parskip}
\usepackage{tocloft}
\usepackage{lastpage}
\usepackage{cite}

\geometry{top=2.5cm, bottom=2.5cm, left=3cm, right=2.5cm}
\setstretch{1.15}
\widowpenalty=10000
\clubpenalty=10000

% ------ Code style ------
\lstdefinestyle{aegiscode}{
  backgroundcolor=\color{gray!10},
  basicstyle=\ttfamily\small,
  breaklines=true,
  frame=single,
  rulecolor=\color{gray!50},
  numbers=left,
  numberstyle=\tiny\color{gray},
  keywordstyle=\color{blue!70},
  commentstyle=\color{green!50!black},
  stringstyle=\color{red!60},
  tabsize=2,
  captionpos=b
}

% ------ Header / Footer ------
\pagestyle{fancy}
\fancyhf{}
\lhead{\small Aegis OS --- AID 325 Final Report}
\rhead{\small E-JUST, May 2026}
\cfoot{Page \thepage\ of \pageref{LastPage}}
\renewcommand{\headrulewidth}{0.4pt}

\hypersetup{
  colorlinks=true,
  linkcolor=blue!60!black,
  citecolor=green!50!black,
  urlcolor=blue!60!black,
  pdftitle={Aegis OS -- AID 325 Final Report},
  pdfauthor={Omar Dorgham, Mariam El-Kholey}
}

% ============================================================
\begin{document}
% ============================================================

% ---- COVER PAGE -------------------------------------------
\begin{titlepage}
  \thispagestyle{empty}
  \centering
  \vspace*{1.5cm}

  {\LARGE\bfseries Aegis OS: An Integrated AI Accountability and\\[6pt]
  Blockchain-Enforced Cyber Insurance\\[6pt]
  Verification Platform\par}

  \vspace{0.6cm}
  {\large\itshape Final Project Report --- AID 325: Blockchain \& Distributed Ledgers\par}

  \vspace{1.5cm}
  \rule{0.6\linewidth}{0.4pt}\\[0.6cm]
  {\large
    \textbf{Omar Dorgham} \quad ID: 320230044\\[4pt]
    \textbf{Mariam El-Kholey} \quad ID: 320230093
  }\\[0.6cm]
  \rule{0.6\linewidth}{0.4pt}

  \vspace{1.2cm}
  {\normalsize
    \textbf{Supervisor:} Dr.\ Muhammad Hataba\\
    Assistant Professor, Computer Science \& Engineering Department\\[6pt]
    \textbf{Co-Supervisor:} Eng.\ Salma Walid
  }

  \vspace{1.2cm}
  {\normalsize
    Egypt-Japan University of Science and Technology (E-JUST)\\
    Faculty of Engineering\\
    Department of Computer Science \& Information Technology
  }

  \vfill
  {\large May 18, 2026}
\end{titlepage}

\clearpage
\pagenumbering{roman}

% ---- ABSTRACT ---------------------------------------------
\begin{abstract}
\noindent
The convergence of autonomous artificial intelligence and cyber insurance
creates two structurally linked accountability crises. First, high-risk AI
systems that take real-world actions produce reasoning logs stored in
mutable relational databases; any database administrator can silently alter
or delete these records, making post-hoc accountability impossible. Second,
the global cyber insurance market---reaching \$16.3 billion in gross
premiums in 2025---suffers from a bilateral evidence problem: insureds
misrepresent their pre-breach security posture, while attackers deliberately
destroy log servers and backup systems before filing claims, leaving
adjudicators with no neutral ground truth. Both crises share one structural
root cause: a single party controls the evidence they are accused of
falsifying. Blockchain eliminates this assumption through cryptographic
immutability.

\emph{Aegis OS} is a working prototype that addresses both problems through
a unified hybrid architecture. The system deploys seven Solidity smart
contracts on a Hardhat-managed Ethereum-compatible network:
\texttt{AuditRegistry}, \texttt{PostureRegistry}, \texttt{PolicyEngine},
\texttt{ClaimsProcessor}, \texttt{CyberToken} (ERC-20), \texttt{Governance},
and \texttt{DailyReportRegistry}. Three Python-based AI agents powered by
Google Gemini via LangChain---\texttt{SecurityAgent},
\texttt{AnomalyDetector}, and \texttt{FraudAnalyzer}---operate as the
intelligent layer, each producing Pydantic-validated structured decisions
that are cryptographically hashed and anchored on-chain. An asynchronous
Merkle notarization pipeline batches agent events every 60 seconds,
reducing per-event blockchain transaction costs by a factor of
approximately 1,530$\times$ compared to synchronous per-event commitment.
Full reasoning text is preserved in a local SQLite database, while only
Merkle roots are committed on-chain, balancing gas efficiency with
forensic completeness.

The key architectural innovation is the integration of the AI audit
subsystem with the parametric insurance subsystem: \texttt{ClaimsProcessor.processClaim()}
cross-queries \texttt{AuditRegistry} to count AI defensive actions in the
seven days preceding a reported breach and applies a ten-point Security
Hygiene Score (SHS) bonus when more than five defensive actions were logged.
This creates a verifiable, tamper-proof link between pre-breach AI
defensive activity and insurance claim outcomes.

The system demonstrates compliance with EU AI Act Regulation (EU)
2024/1689, Article 12, which mandates tamper-resistant automatic event
logging for high-risk AI systems by August 2, 2026. A full-stack React
interface with MetaMask wallet integration, ERC-20 token payouts, and
adversarial reentrancy testing via \texttt{MaliciousAttacker.sol} complete
the prototype. Future work includes Sepolia testnet deployment, formal
contract verification with Certora Prover, and zero-knowledge proof
integration to enable privacy-preserving AI reasoning audits.
\end{abstract}

\clearpage
\tableofcontents
\clearpage

\pagenumbering{arabic}
\setcounter{page}{1}

% ============================================================
\section{Introduction}
\label{sec:intro}
% ============================================================

\subsection{Context and Motivation}
\label{sec:motivation}

The regulatory landscape for autonomous AI systems is undergoing a
structural shift. EU AI Act (Regulation (EU) 2024/1689), Article 12,
mandates that providers of high-risk AI systems implement
``automatic recording of events over the lifetime of the system''
using ``tamper-resistant'' logging mechanisms, effective August 2, 2026.
Non-compliance carries penalties of up to \texteuro35 million or 7\%
of global annual turnover, whichever is higher \cite{euaiact2024}.
No currently deployed commercial security operations platform satisfies
this requirement using cryptographically verifiable, externally auditable
logging infrastructure.

Simultaneously, the cyber insurance industry faces an evidence crisis
of its own. The global cyber insurance market reached \$16.3 billion in
gross written premiums in 2025, projected to grow at a 27\% compound
annual growth rate through 2030, reaching approximately \$53 billion
\cite{cybermarket2024}. Despite this scale, claims resolution averages
six to eighteen months due to three structural failures: (1) insureds
misrepresent their pre-breach security posture by retroactively claiming
controls that did not exist; (2) ransomware attackers deliberately target
log servers and backup infrastructure before breach notification, destroying
the evidence base; and (3) both parties distrust evidence produced or
controlled by the opposing party.

Both problems share a single structural root cause: a trusted---and
potentially adversarial---party controls the evidence. Immutable blockchain
ledgers eliminate this assumption by design. Once data is committed to a
distributed ledger under consensus, it cannot be altered without detection.
\emph{Aegis OS} exploits this property to build the first unified
prototype that addresses both the AI accountability gap and the cyber
insurance evidence problem in a single integrated system.

\subsection{Problem Statement}
\label{sec:problem}

The system addresses three formally distinct but structurally linked
problems:

\begin{description}[leftmargin=2em]
  \item[\textbf{P1 (AI Accountability):}] Given that an AI agent takes
  autonomous actions with real-world consequences, how can a third party
  verify---after the fact---exactly what the agent decided, what context
  it had, and what reasoning it applied, without trusting the log source?

  \item[\textbf{P2 (Insurance Adjudication):}] Given that a cyber
  insurance claim involves two adversarial parties and that the breach
  may have destroyed the pre-breach evidence, how can claim adjudication
  be made objective, automated, and fraud-resistant?

  \item[\textbf{P3 (Integration):}] How can pre-breach AI defensive
  activity be incorporated as evidence in insurance claim adjudication,
  such that organizations that actively defended are rewarded
  proportionally and verifiably?
\end{description}

\subsection{Research Objectives}
\label{sec:objectives}

\begin{enumerate}[label=\textbf{O\arabic*.}]
  \item Design and implement a tamper-proof AI audit trail using
        cryptographic SHA-256 hashing, binary Merkle trees, and
        blockchain anchoring via \texttt{AuditRegistry.sol}.
  \item Implement automated cyber insurance claims processing using
        on-chain \texttt{ClaimsProcessor.sol} logic with no human
        intermediary in the verdict path.
  \item Integrate AI defensive action evidence into insurance payout
        calculations through cross-contract querying between
        \texttt{ClaimsProcessor} and \texttt{AuditRegistry}.
  \item Demonstrate EU AI Act Article 12 compliance through technical
        implementation of automatic, tamper-resistant, reconstructible
        event logging.
  \item Validate the system through Hardhat unit tests, integration
        testing, and adversarial security testing using
        \texttt{MaliciousAttacker.sol} with OpenZeppelin
        \texttt{ReentrancyGuard}.
  \item Demonstrate gas efficiency through Merkle batching and provide
        a forensic gas cost analysis across all primary contract functions.
\end{enumerate}

\subsection{Contributions}
\label{sec:contributions}

\begin{itemize}
  \item A working prototype that integrates AI audit accountability with
        parametric insurance verdict logic in a single verifiable on-chain
        system---no equivalent commercial product exists.
  \item An asynchronous Merkle notarization pipeline
        (\texttt{notarizer.py}) with sequential event ID integrity,
        atomic persistence via \texttt{os.replace}, and automatic gap
        detection that makes event dropping forensically detectable.
  \item A dual-storage architecture that separates full-fidelity
        reasoning data (SQLite, JSON batch logs) from cryptographic
        anchors (blockchain Merkle roots), achieving gas efficiency
        without sacrificing auditability.
  \item An adversarial test harness (\texttt{MaliciousAttacker.sol})
        that deploys a contract simulating a recursive reentrancy attack,
        verifying that the two-layer defence (CEI pattern plus
        \texttt{nonReentrant}) holds under realistic attack conditions.
\end{itemize}

\subsection{Report Organization}
\label{sec:organization}

Section~\ref{sec:background} surveys the technical background on AI
accountability, cyber insurance evidence problems, Merkle tree audit
infrastructure, and parametric insurance smart contracts.
Section~\ref{sec:architecture} presents the five-layer system
architecture, dual-storage design rationale, and the 14-step runtime
data flow. Section~\ref{sec:contracts} documents all eight smart
contracts with function-level detail. Section~\ref{sec:backend} covers
the Python backend including the SQLite schema, AI agent architecture,
notarizer, and FastAPI endpoints. Section~\ref{sec:frontend} documents
the React frontend and MetaMask integration. Section~\ref{sec:security}
provides a comprehensive security analysis including reentrancy
protection, cryptographic guarantees, and a formal threat model.
Section~\ref{sec:testing} presents the full testing and validation
results. Section~\ref{sec:compliance} addresses regulatory compliance
and market positioning. Sections~\ref{sec:future} and~\ref{sec:conclusion}
present future work and conclusions, respectively.

% ============================================================
\section{Background and Related Work}
\label{sec:background}
% ============================================================

\subsection{The AI Accountability Problem}
\label{sec:ai-accountability}

Traditional database-backed logging systems fail the requirements of
EU AI Act Article 12 on four structural grounds. First, any user holding
\texttt{DELETE} or \texttt{UPDATE} privileges on the log table can silently
alter or remove records; SQL does not provide cryptographic integrity
guarantees by design. Second, log timestamps are generated by the
application server's system clock, which the operating entity controls
and can manipulate. Third, log integrity cannot be proven to any external
party without trusting the entity that operates the database---precisely
the party whose conduct is under scrutiny. Fourth, there is no
cryptographic binding between a log entry's content and its existence in
the database; a row can be replaced with fabricated content and the
database remains internally consistent.

Article 12 of Regulation (EU) 2024/1689 specifies three requirements
that directly address these failures: ``automatic recording of events
over the lifetime of the system,'' the ability to ``enable the
reconstruction of how the system produced a specific output,'' and
``tamper-resistant logging mechanisms'' \cite{euaiact2024}. No
commercially deployed security operations platform (Palo Alto XSOAR,
CrowdStrike Falcon, Splunk SIEM) currently satisfies the
\emph{tamper-resistant} requirement through cryptographically verifiable,
externally auditable blockchain anchoring.

\subsection{The Cyber Insurance Evidence Problem}
\label{sec:insurance-problem}

Cyber insurance claims adjudication suffers from a structural bilateral
fraud problem. The insured party has incentives to misrepresent pre-breach
security posture---claiming that multi-factor authentication, offsite
backups, and vulnerability patching were in place when they were not.
The insurer has incentives to dispute claims by citing policy exclusions
or negligence clauses. Ransomware operators are known to specifically
target log aggregation servers and backup infrastructure prior to
detonating their payload, destroying the evidentiary base before the
claim is even filed \cite{ransomware2023}.

Table~\ref{tab:db-vs-blockchain} compares the structural properties of
traditional database logs versus blockchain-anchored records.

\begin{table}[ht]
\centering
\caption{Traditional Database vs.\ Blockchain Ledger for Evidence Storage}
\label{tab:db-vs-blockchain}
\begin{tabular}{@{}lll@{}}
\toprule
\rowcolor{gray!15}
\textbf{Property} & \textbf{Traditional Database} & \textbf{Blockchain Ledger} \\
\midrule
Control         & Single entity (defendant)       & No single party \\
Modification    & Possible by administrator       & Cryptographically impossible \\
Timestamps      & Company server (mutable)        & Block headers (immutable) \\
Integrity proof & Cannot be proven externally     & Merkle proof, verifiable by any party \\
Admissibility   & Disputed by opposing party      & Neutral, deterministic witness \\
\bottomrule
\end{tabular}
\end{table}

\subsection{Merkle Trees as Audit Infrastructure}
\label{sec:merkle}

A Merkle tree \cite{merkle1987} is a binary hash tree in which every
leaf node contains the hash of a data block and every internal node
contains the hash of its two children. The root of the tree is a single
32-byte value that cryptographically commits to the entirety of the
underlying data. The key operational property is $O(\log n)$ proof size:
to prove that a specific event belongs to a batch of $n$ events, an
auditor needs only $\log_2 n$ sibling hashes, not the full dataset.

A known vulnerability is the \emph{second-preimage attack}: without
domain separation, an attacker can present an internal node hash as a
leaf hash, and the Merkle proof will verify correctly despite the data
being fabricated. The standard mitigation \cite{laurie2012} uses a
one-byte domain prefix: leaves are hashed as
$\text{SHA256}(\texttt{0x00} \| \text{data})$ and internal nodes as
$\text{SHA256}(\texttt{0x01} \| L \| R)$. Aegis OS implements this
exactly in \texttt{notarizer.py} lines 100--102 and 148--149:
\begin{lstlisting}[style=aegiscode, language=Python,
  caption={Domain-separated Merkle hashing in notarizer.py}]
# Leaf: SHA256(0x00 + payload)
event_hash = "0x" + hashlib.sha256(b"\x00" + payload.encode()).hexdigest()

# Node: SHA256(0x01 + left + right)
next_level.append(hashlib.sha256(b"\x01" + left + right).digest())
\end{lstlisting}

An auditor wishing to verify an event retrieves the Merkle root from the
immutable blockchain via \texttt{AuditRegistry.getEntry()}, then
recomputes the Merkle path from the off-chain batch JSON file and
confirms the path resolves to the on-chain root. Any tampering with the
off-chain JSON invalidates the proof.

\subsection{Parametric Insurance and Smart Contracts}
\label{sec:parametric}

Traditional insurance operates on discretionary human judgment:
adjusters investigate claims, negotiate settlements, and make subjective
decisions that can be disputed for years. Parametric insurance replaces
this with objective, pre-agreed trigger conditions: if a verifiable
parameter (e.g., windspeed at a weather station, or a cryptographically
attested security score) crosses a threshold, a payout is automatically
triggered---no discretion, no negotiation \cite{parametric2022}.

Ethereum smart contracts are structurally ideal parametric insurance
engines because their logic is immutable once deployed, execution is
deterministic across all nodes, and they can transfer value
(ERC-20 tokens) without human intermediaries. \texttt{ClaimsProcessor.sol}
implements this architecture: the payout formula is encoded in
Solidity bytecode, deployed to the chain, and cannot be changed by
either party after a policy is activated.

\subsection{Comparison with Existing Systems}
\label{sec:comparison}

Table~\ref{tab:system-comparison} positions Aegis OS against existing
commercial and open-source platforms across the three capability
dimensions the system addresses.

\begin{table}[ht]
\centering
\caption{System Comparison: AI Audit, Blockchain Anchoring, Insurance Integration}
\label{tab:system-comparison}
\begin{tabular}{@{}llll@{}}
\toprule
\rowcolor{gray!15}
\textbf{System} & \textbf{AI Audit} & \textbf{Blockchain} & \textbf{Insurance Integration} \\
\midrule
Palo Alto XSOAR   & Yes     & No  & No \\
CrowdStrike Falcon & Yes    & No  & No \\
Traditional SIEM  & Partial & No  & No \\
\textbf{Aegis OS} & \textbf{Yes} & \textbf{Yes} & \textbf{Yes (automated)} \\
\bottomrule
\end{tabular}
\end{table}

No existing commercial product provides all three capabilities in a
single verifiable system. Aegis OS is a working prototype of this
previously unimplemented architecture.


% ============================================================
\section{System Architecture}
\label{sec:architecture}
% ============================================================

\subsection{Architectural Philosophy}
\label{sec:arch-philosophy}

Aegis OS is governed by two complementary trust assumptions derived from
the evidence problems stated in Section~\ref{sec:problem}.
First, no party should be trusted to store evidence they are accused of
falsifying; blockchain enforces this by making stored Merkle roots
cryptographically immutable. Second, expensive operations---storing full
reasoning text, all event metadata, and query-friendly relational
data---should not happen on-chain, where every byte incurs gas cost.
This produces a \emph{dual-storage architecture}:

\begin{itemize}
  \item \textbf{Off-chain (SQLite + JSON):} Full-fidelity data including
        complete LLM reasoning text, action descriptions, agent metadata,
        and batch event lists. Fast, cheap, queryable by the frontend via
        REST API.
  \item \textbf{On-chain (Solidity contracts):} Cryptographic anchors
        only---Merkle roots, posture snapshot hashes, compliance verdicts,
        and ERC-20 token balances. Immutable, publicly auditable, gas-
        efficient.
\end{itemize}

An external auditor can verify any off-chain record against its on-chain
anchor using a Merkle proof. Neither layer can be tampered with
independently without the discrepancy becoming detectable.

\subsection{Five-Layer Stack}
\label{sec:layers}

The system is organized into five horizontal layers, each with a specific
responsibility and a justified technology choice.

\begin{description}[leftmargin=2em]
  \item[\textbf{Layer 1 --- Frontend (React + Vite):}]
  React's component model maps directly to role-based dashboard views
  (Admin and Client roles render entirely different page sets). Vite
  provides sub-second hot-module reload during development and an
  optimized production bundle. The frontend reads on-chain state through
  \texttt{ethers.js} \texttt{BrowserProvider} and off-chain data through
  the FastAPI REST layer.

  \item[\textbf{Layer 2 --- Backend AI/API (FastAPI + Agents):}]
  FastAPI is async-native, matches LangChain's \texttt{async} chain
  interface without bridging overhead, and auto-generates OpenAPI
  documentation at \texttt{/docs}. Pydantic models used by FastAPI are
  identical to the structured output models used by LangChain agents,
  eliminating a translation layer. Three single-responsibility agents
  (\texttt{SecurityAgent}, \texttt{AnomalyDetector}, \texttt{FraudAnalyzer})
  handle distinct temporal scopes and output types.

  \item[\textbf{Layer 3 --- Integrity/Notarization (notarizer.py):}]
  Notarization is a background process that must not block AI agent
  response latency. An \texttt{asyncio.Queue} decouples agent event
  production (sub-millisecond enqueue) from blockchain submission
  (network-bound). \texttt{APScheduler AsyncIOScheduler} fires
  \texttt{run\_batch()} every 60 seconds on the same event loop as the
  FastAPI application, avoiding thread-safety issues.

  \item[\textbf{Layer 4 --- Blockchain (Hardhat + Solidity + OpenZeppelin):}]
  Hardhat is chosen over Truffle for its native gas reporting plugin,
  faster recompile cycle, and superior TypeScript ecosystem integration.
  OpenZeppelin v5 contracts provide audited, battle-tested base
  implementations for \texttt{ERC20}, \texttt{ReentrancyGuard}, and
  \texttt{Ownable}, reducing the attack surface compared to custom
  implementations. Solidity 0.8.28 is specified in \texttt{hardhat.config.js}.

  \item[\textbf{Layer 5 --- Persistence (SQLite + JSON batch logs):}]
  SQLite is zero-configuration, single-file, and appropriate for
  prototype scale. The schema (defined in \texttt{database.py}) can be
  migrated to PostgreSQL for production without changing any application
  logic. JSON batch logs in \texttt{audit\_log/} serve as the primary
  audit artifact: they are human-readable, contain full event data, and
  can be verified against on-chain Merkle roots without a running database.
\end{description}

\subsection{Runtime Data Flow}
\label{sec:dataflow}

The following 14-step sequence traces the complete lifecycle from a
frontend user action to a finalized on-chain record:

\begin{enumerate}
  \item The React frontend sends an HTTP POST request to a FastAPI
        endpoint (e.g., \texttt{/api/analyze-alert}).
  \item FastAPI deserializes the request body via a Pydantic model and
        routes it to the appropriate AI agent's handler method.
  \item The AI agent invokes the Gemini LLM via a LangChain
        \texttt{PromptTemplate | LLM | structured\_output} chain
        using \texttt{chain.ainvoke()}.
  \item The LLM response is deserialized into a Pydantic output model
        (\texttt{SecurityDecision}, \texttt{AnomalyReport}, or
        \texttt{FraudAnalysis}), enforcing field types and value bounds.
  \item The agent persists the full decision (instruction, context,
        reasoning, action, risk level) to the \texttt{ai\_decisions}
        SQLite table via \texttt{database.insert\_ai\_decision()}.
  \item For high-risk events (risk level $\geq 2$), the agent computes
        four SHA-256 hashes (instruction, context, reasoning, action)
        and enqueues an event dict to the shared \texttt{asyncio.Queue}.
  \item \texttt{APScheduler} fires \texttt{notarizer.run\_batch()} every
        60 seconds on the shared event loop.
  \item The notarizer drains the queue, assigns sequential IDs via an
        \texttt{asyncio.Lock}-protected counter, and checks ID contiguity
        to detect dropped events.
  \item A domain-separated Merkle tree is constructed from the batch:
        leaf hashes use the \texttt{0x00} prefix; node hashes use \texttt{0x01}.
  \item A single blockchain transaction submits the Merkle root to
        \texttt{AuditRegistry.logAction()} (or to
        \texttt{DailyReportRegistry.submitDailyReport()} for daily batches).
  \item The batch JSON is atomically written to \texttt{audit\_log/}
        using \texttt{os.replace()} to prevent partial-write corruption.
  \item The frontend polls \texttt{/api/activity-feed} every 5 seconds
        and reads on-chain state directly via \texttt{ethers.js} contract calls.
  \item For insurance: the client calls \texttt{ClaimsProcessor.fileClaim()};
        the backend \texttt{FraudAnalyzer} calls \texttt{recordFraudScore()};
        the backend then calls \texttt{processClaim()}, which internally
        cross-queries both \texttt{PostureRegistry} and \texttt{AuditRegistry}.
  \item On verdict \texttt{APPROVED} or \texttt{PARTIAL},
        \texttt{ClaimsProcessor} transfers \texttt{CyberToken} ERC-20
        tokens directly to the claimant's wallet.
\end{enumerate}

\subsection{AI Accountability Sub-System}
\label{sec:ai-subsystem}

\subsubsection{The Five-Hash Pipeline}
\label{sec:fivehash}

Every AI agent decision is decomposed into four independently hashed
components before blockchain commitment. The design rationale for
separating the hashes is forensic isolation: if all components were
combined into one hash, an auditor could not determine \emph{which}
component changed. With four separate hashes, it becomes possible to
state precisely: ``the instruction and context hashes match the
reference, but the reasoning hash differs---the agent produced different
reasoning for the same input.''

\begin{description}[leftmargin=2em]
  \item[\textbf{Instruction Hash} \texttt{SHA256(prompt\_template)}:]
  Proves exactly what the agent was instructed to do. Stored separately
  so instruction changes are detectable independently of context changes.
  \item[\textbf{Context Hash} \texttt{SHA256(system\_context)}:]
  Proves what information the agent had access to at decision time.
  Critical for proving the agent did not have access to information it
  claims it did not have.
  \item[\textbf{Reasoning Hash} \texttt{SHA256(decision.reasoning)}:]
  Proves the logical chain from context to action. Stored separately so
  reasoning can be verified without exposing full content to the insurer.
  \item[\textbf{Action Hash} \texttt{SHA256(decision.action)}:]
  Proves what was actually executed, not merely decided.
\end{description}

In \texttt{security\_agent.py} lines 60--78, the agent constructs all
four hashes and enqueues them immediately after the LLM responds:
\begin{lstlisting}[style=aegiscode, language=Python,
  caption={Four-hash pipeline in SecurityAgent.handle\_alert()}]
ih = hasher.sha256(self.prompt.template)      # instruction
ch = hasher.sha256(system_context)             # context
rh = hasher.sha256(decision_obj.reasoning)     # reasoning
ah = hasher.sha256(decision_obj.action)        # action

blockchain_id = blockchain.log_ai_action(ih, ch, rh, ah, decision_obj.risk_level)
event = {"event_id": blockchain_id, "instruction_hash": ih,
         "context_hash": ch, "reasoning_hash": rh, "action_hash": ah, ...}
await queue.put(event)
\end{lstlisting}

\subsubsection{The Async Notarizer}
\label{sec:notarizer-arch}

The notarizer's 60-second batch interval is a deliberate economic
optimization. At 100 alerts per minute, synchronous per-event commitment
would cost approximately \$765/minute at mainnet gas prices
(\$7.65/transaction $\times$ 100). Batching 100 events into one Merkle
root costs approximately \$0.50/minute---a factor of 1,530$\times$
reduction with identical cryptographic security guarantees.

The \texttt{max\_instances=1} and \texttt{coalesce=True} scheduler
settings prevent queue double-drain: if the blockchain is slow and a
scheduled fire is missed, \texttt{coalesce=True} executes one catch-up
run rather than all missed runs. \texttt{max\_instances=1} prevents a
second \texttt{run\_batch()} from starting while the first is still
submitting a transaction.

\subsection{Cyber Insurance Sub-System}
\label{sec:insurance-subsystem}

\subsubsection{Daily Posture Snapshots}

\texttt{PostureRegistry.sol} stores daily cryptographic snapshots of
each company's security posture. Each \texttt{PostureSnapshot} struct
contains a \texttt{merkleRoot} (a 32-byte hash representing the composite
of firewall ruleset hash, backup snapshot hash, CVE patch coverage hash,
and admin ACL hash), a Unix \texttt{timestamp}, the submitting
\texttt{company} address, and a \texttt{complianceScore} (0--100).
Daily---rather than real-time---snapshots are used because they create
a timestamped baseline: the question ``what was the security posture
at 2pm on May 15?'' is answered by ``the last verified snapshot before
that time,'' preventing infinite on-chain data growth.

\subsubsection{PolicyEngine Design}

The insurance policy is encoded on-chain in \texttt{PolicyEngine.sol}
rather than in a PDF or off-chain database. This design decision is
critical: if the policy document is stored off-chain, the insurer can
claim the terms were different after a breach occurs. On-chain policy
encoding cryptographically binds both parties to identical, immutable
terms from the moment the policy is activated---before any breach occurs.

\subsubsection{Automated Claims Processing}

\texttt{ClaimsProcessor.processClaim()} implements a fully automated
parametric verdict with the following decision logic derived directly
from the contract code (lines 139--187):

\begin{equation}
\label{eq:payout}
\text{PayoutAmount} = \text{claimedAmount} \times \frac{\text{SHS}}{100}
\end{equation}

where SHS is the 90-day rolling Security Hygiene Score from
\texttt{PostureRegistry}, subject to:
\begin{itemize}
  \item \textbf{AI Defense Boost:} if \texttt{auditRegistry.getActionCountByTimeRange}
        returns $> 5$ actions in the 7-day window before breach,
        SHS is boosted by 10 points (capped at 100).
  \item \textbf{APPROVED} (100\% payout): SHS $\geq 90$ AND
        \texttt{fraudScore} $< 30$.
  \item \textbf{PARTIAL} (SHS\% payout): SHS $\geq 50$.
  \item \textbf{DENIED} (0\% payout): SHS $< 50$ OR
        \texttt{fraudScore} $> 75$.
\end{itemize}

\subsection{The Integration Layer}
\label{sec:integration}

\subsubsection{Architectural Innovation}

Module A (\texttt{AuditRegistry}) alone answers: ``What did the AI
decide, and can we prove it?'' Module B (\texttt{ClaimsProcessor})
alone answers: ``Was the company's security posture compliant when the
breach occurred?'' Their integration answers the question no existing
commercial system can: ``Given that AI was actively defending---and all
its actions are provably recorded---what does the evidence say about
this breach, and what payout does the evidence justify?''

\subsubsection{Cross-Contract Evidence Query}

The integration mechanism is implemented in
\texttt{ClaimsProcessor.processClaim()} lines 151--158:

\begin{lstlisting}[style=aegiscode, language=Solidity,
  caption={Cross-contract AI evidence query in processClaim()}]
uint256 startTime = claim.breachTimestamp > 7 days
    ? claim.breachTimestamp - 7 days : 0;
uint256 aiActions = auditRegistry.getActionCountByTimeRange(
    startTime, claim.breachTimestamp);

if (aiActions > 5) {
    uint16 boostedSHS = uint16(shs) + 10;
    shs = boostedSHS > 100 ? 100 : uint8(boostedSHS);
}
\end{lstlisting}

The 7-day window is chosen to capture active defensive posture
immediately preceding the breach without penalizing companies whose
AI agent was inactive during unrelated quiet periods.

% ============================================================
\section{Smart Contract Implementation}
\label{sec:contracts}
% ============================================================

\subsection{AuditRegistry.sol}
\label{sec:audit-registry}

\texttt{AuditRegistry} is the immutable append-only ledger for AI
agent decisions. Its core data structure is a public dynamic array
\texttt{AuditEntry[] public entries}, where each entry stores four
32-byte hashes (\texttt{instructionHash}, \texttt{contextHash},
\texttt{reasoningHash}, \texttt{actionHash}), a
\texttt{RiskLevel} enum value, and a \texttt{block.timestamp}.
The design is intentionally append-only: no \texttt{delete} or
\texttt{update} functions exist, making the log structurally immutable
at the Solidity level.

\texttt{logAction()} is \texttt{external} with no access modifier,
meaning any caller (including the backend system) can append entries.
This is appropriate because the security guarantee comes not from
access restriction but from cryptographic commitment: once appended,
an entry cannot be modified without changing the block hash.
The function returns \texttt{entryId} (the array index) and emits
\texttt{ActionLogged(entryId, actionHash, risk)}.

\texttt{getActionCountByTimeRange(startTimestamp, endTimestamp)}
iterates all entries and counts those whose \texttt{timestamp} falls
within the range. This linear scan is the primary gas cost driver for
the cross-contract call in \texttt{ClaimsProcessor}.

\subsection{Governance.sol}
\label{sec:governance}

\texttt{Governance} implements a human-oversight approval mechanism
for HIGH and CRITICAL risk AI actions. It holds an
\texttt{immutable} reference to \texttt{AuditRegistry} and two
state mappings: \texttt{pendingApprovals} and \texttt{approvedActions},
both \texttt{mapping(uint256 => bool)}.

\texttt{requestApproval(\_entryId)} fetches the entry from
\texttt{AuditRegistry.getEntry()} and requires
\texttt{uint8(entry.riskLevel) >= 2} (HIGH or CRITICAL). This prevents
governance queue flooding with low-risk routine actions. The
\texttt{onlySigner} modifier restricts \texttt{approve()} and
\texttt{reject()} to addresses explicitly added by the \texttt{onlyOwner}-
protected \texttt{addSigner()} function. The owner cannot be removed
(\texttt{require(\_signer != owner, "Cannot remove owner")}),
preventing a governance lockout attack.

\subsection{PostureRegistry.sol}
\label{sec:posture-registry}

\texttt{PostureRegistry} stores daily security posture evidence.
The \texttt{PostureSnapshot} struct (merkleRoot, timestamp, company,
complianceScore) is stored in a private
\texttt{mapping(address => PostureSnapshot[])} to prevent direct
external array access. \texttt{recordSnapshot(\_merkleRoot, \_score)}
requires \texttt{\_score <= 100} to prevent invalid input, then
appends a new snapshot using \texttt{msg.sender} as the company
address, emitting \texttt{SnapshotRecorded}.

\texttt{getSecurityHygieneScore(\_company)} computes a rolling average
over the most recent \texttt{min(count, 90)} snapshots, implementing
the 90-day window. The \texttt{uint8} return type bounds the score
to 0--255 by the EVM, but the \texttt{<= 100} input validation and
averaging logic keep it within 0--100.

\subsection{PolicyEngine.sol}
\label{sec:policy-engine}

\texttt{PolicyEngine} stores insurance coverage rules on-chain as
\texttt{PolicyRule} structs (ruleName, required, maxAgeSeconds, weight).
Rules are stored in \texttt{mapping(address => PolicyRule[])} per
company. The \texttt{onlyInsurer} modifier restricts \texttt{setPolicy()}
to the \texttt{immutable} insurer address set in the constructor,
preventing companies from modifying their own coverage terms.
\texttt{setPolicy()} deletes the existing rule array before pushing
new rules, ensuring a clean update without array length inconsistencies.
\texttt{isPolicyActive(\_company)} returns the boolean from the
corresponding mapping, used as a gating check in
\texttt{ClaimsProcessor.fileClaim()}.

\subsection{ClaimsProcessor.sol}
\label{sec:claims-processor}

\texttt{ClaimsProcessor} is the most complex contract and the
integration centrepiece. It inherits OpenZeppelin
\texttt{ReentrancyGuard} and holds four \texttt{immutable} interface
references: \texttt{IPostureRegistry}, \texttt{IPolicyEngine},
\texttt{IAuditRegistry}, and \texttt{IERC20} (\texttt{CyberToken}).
The \texttt{immutable} keyword prevents SLOAD operations on every
access (saving approximately 2,100 gas per read compared to regular
\texttt{storage} variables).

The \texttt{Claim} struct packs \texttt{verdict} (enum, 1 byte),
\texttt{payoutPercentage} (uint8, 1 byte), and \texttt{fraudScore}
(uint8, 1 byte) into a single 32-byte EVM slot, saving approximately
40,000 gas compared to separate storage slots.

\texttt{processClaim()} implements the Checks-Effects-Interactions
pattern rigorously:

\begin{lstlisting}[style=aegiscode, language=Solidity,
  caption={CEI pattern in ClaimsProcessor.processClaim() (condensed)}]
function processClaim(uint256 _claimId)
    external onlyBackend nonReentrant {

  // 1. CHECKS
  require(_claimId < claimCount, "Claim does not exist");
  require(claim.verdict == Verdict.PENDING, "Claim already processed");
  require(claim.fraudReportHash != bytes32(0),
      "AI Fraud score must be recorded first");

  // 2. EFFECTS (all state changes before external calls)
  claim.verdict = Verdict.APPROVED; // or PARTIAL / DENIED
  claim.payoutPercentage = ...;
  claim.processedAt = block.timestamp;
  emit ClaimProcessed(_claimId, claim.verdict, payoutAmount);

  // 3. INTERACTIONS (external call last)
  if (payoutAmount > 0) {
    require(cyberToken.transfer(claim.company, payoutAmount),
        "Token transfer failed");
  }
}
\end{lstlisting}

The \texttt{nonReentrant} modifier (from OpenZeppelin
\texttt{ReentrancyGuard}) sets a mutex lock before the function body
executes and releases it after. Any reentrant call finds the lock set
and reverts with \texttt{ReentrancyGuardReentrantCall}. The CEI pattern
provides a second layer: even if \texttt{nonReentrant} were removed,
updating \texttt{claim.verdict} before the token transfer means any
reentrant call would fail the \texttt{require(claim.verdict == PENDING)}
check.

\subsection{CyberToken.sol}
\label{sec:cyber-token}

\texttt{CyberToken} inherits OpenZeppelin \texttt{ERC20} and
\texttt{Ownable}. The constructor calls \texttt{\_mint(msg.sender,
initialSupply)}, minting all tokens to the deployer who then transfers
50,000 CIT to \texttt{ClaimsProcessor} via \texttt{deploy.js}.
The \texttt{onlyOwner}-restricted \texttt{mint()} function allows
future supply expansion. ERC-20 is chosen over a custom token
implementation for three reasons: MetaMask and all EVM wallets natively
display ERC-20 balances, the OpenZeppelin base contract is formally
audited, and any exchange or DeFi protocol can integrate the token
without custom adapters.

\subsection{DailyReportRegistry.sol}
\label{sec:daily-report}

\texttt{DailyReportRegistry} stores daily aggregate summaries of AI
activity. The \texttt{DailyReport} struct contains \texttt{reportHash}
(bytes32), \texttt{timestamp} (uint256), \texttt{totalActions}
(uint256), \texttt{avgRiskLevel} (uint8), and \texttt{submitted} (bool).
Reports are stored in \texttt{mapping(uint256 => DailyReport)} indexed
by day number. \texttt{submitDailyReport()} is restricted by
\texttt{onlyBackend}, stores the report, emits
\texttt{DailyReportSubmitted}, and increments \texttt{currentDay}.
\texttt{getReport()} requires \texttt{submitted == true} to prevent
reading uninitialised slots.

\subsection{MaliciousAttacker.sol}
\label{sec:malicious-attacker}

\texttt{MaliciousAttacker} is an adversarial test contract that
simulates a reentrancy attack against \texttt{ClaimsProcessor}.
Its \texttt{deployTarget()} function deploys a new
\texttt{ClaimsProcessor} instance with the attacker contract as
\texttt{backendSystem} \emph{and} as the \texttt{auditRegistry}
(since \texttt{MaliciousAttacker} also implements
\texttt{recordFraudScore}). This gives the attacker full control over
both the backend role and the external call target.

\texttt{attack(\_claimId)} sets \texttt{attackClaimId} and calls
\texttt{targetContract.processClaim(\_claimId)}. During execution,
\texttt{ClaimsProcessor} calls \texttt{auditRegistry.getActionCountByTimeRange()},
which resolves to \texttt{MaliciousAttacker}'s \texttt{fallback()}
function. The \texttt{fallback()} immediately re-enters
\texttt{targetContract.processClaim(attackClaimId)}. Without protection,
this would execute the token transfer twice. With \texttt{nonReentrant},
the re-entrant call reverts with \texttt{ReentrancyGuardReentrantCall},
proving the protection holds under a realistic, contract-level attack
vector---not merely a theoretical one.


% ============================================================
\section{Backend Implementation}
\label{sec:backend}
% ============================================================

\subsection{Database Layer}
\label{sec:database}

The SQLite schema defined in \texttt{database.py} comprises five tables.
\texttt{ai\_decisions} is the primary audit table: columns include \texttt{id}
(PK), \texttt{event\_id} (INT, the blockchain or notarizer sequence ID),
\texttt{agent\_name} (TEXT NOT NULL), \texttt{instruction}, \texttt{context},
\texttt{reasoning}, \texttt{action\_taken} (all TEXT), \texttt{risk\_level}
(TEXT), \texttt{claim\_id} (INT), \texttt{onchain\_entry\_id} (INT), and
\texttt{timestamp} (INT NOT NULL). Each row represents one complete AI
decision cycle, preserving the full reasoning text that is hashed but not
stored on-chain.

\texttt{security\_alerts} links incoming alert text to a decision via
\texttt{decision\_id} (FK to \texttt{ai\_decisions}) and tracks
\texttt{status} (DEFAULT \texttt{'RECEIVED'}).
\texttt{system\_actions} records every categorised action taken
(\texttt{FIREWALL\_BLOCK}, \texttt{QUARANTINE}, \texttt{PATCH\_FLAG}, etc.)
and powers the \texttt{/api/activity-feed} endpoint.
\texttt{claim\_analyses} stores \texttt{FraudAnalyzer} output per claim:
\texttt{fraud\_score}, \texttt{reasoning}, and \texttt{recommendation}
(\texttt{APPROVE}/\texttt{PARTIAL}/\texttt{DENY}).
\texttt{daily\_summaries} aggregates per-day metrics including
\texttt{total\_actions}, \texttt{high\_risk\_count},
\texttt{avg\_risk\_level}, \texttt{report\_hash}, and a
\texttt{notarized} boolean flag.

SQLite is chosen over PostgreSQL for zero-configuration deployment
appropriate to prototype scale. The schema is fully portable: migration
to PostgreSQL requires only changing the connection string.

\subsection{AI Agent Architecture}
\label{sec:agents}

\subsubsection{SecurityAgent}
\label{sec:security-agent}

\texttt{SecurityAgent} handles real-time threat alert analysis.
Its \texttt{SecurityDecision} Pydantic model enforces three fields:
\texttt{reasoning} (str), \texttt{action} (str), and \texttt{risk\_level}
(int, bounded 0--3 via \texttt{ge=0, le=3}). The LangChain chain is
constructed as \texttt{PromptTemplate | ChatGoogleGenerativeAI(temperature=0.2).with\_structured\_output(SecurityDecision)}.
The low temperature (0.2) reduces output variance for security-critical
structured decisions.

When the LLM is unavailable (\texttt{self.llm is None}), the agent
falls back to keyword-based heuristics: alerts containing
\texttt{"CRITICAL"} map to risk level 3, \texttt{"HIGH"} to level 2,
otherwise level 1. This ensures the system degrades gracefully during
LLM outages rather than failing silently.

For \texttt{risk\_level >= 2}, the agent immediately submits hashes to
\texttt{AuditRegistry} via \texttt{blockchain.log\_ai\_action()} and
falls back to \texttt{notarizer.get\_next\_id()} if the blockchain call
returns \texttt{-1}. The event is then enqueued for Merkle batching.

\subsubsection{AnomalyDetector}
\label{sec:anomaly-detector}

\texttt{AnomalyDetector} operates on aggregate daily log data rather
than individual alerts---a different temporal scope requiring a different
prompt structure. Its \texttt{AnomalyReport} model contains
\texttt{anomaly\_detected} (bool), \texttt{reasoning} (str), and
\texttt{recommended\_action} (str). Events are only queued if
\texttt{anomaly\_detected=True}, preventing routine daily summaries
from flooding the notarizer queue. The fallback detects anomalies
based on \texttt{"CRITICAL"} or \texttt{"HIGH"} keywords in the log
JSON string.

\subsubsection{FraudAnalyzer}
\label{sec:fraud-analyzer}

\texttt{FraudAnalyzer} cross-examines insurance claims against historical
posture data. Its \texttt{FraudAnalysis} model enforces
\texttt{fraud\_score} (int, 0--100) and \texttt{report} (str).
Scores below 30 map to \texttt{APPROVE}, 30--74 to \texttt{PARTIAL},
and 75+ to \texttt{DENY}.

Uniquely, \texttt{FraudAnalyzer} also submits a signed transaction
to \texttt{ClaimsProcessor.recordFraudScore()} via
\texttt{build\_transaction()} + \texttt{sign\_transaction(PRIVATE\_KEY)}
+ \texttt{send\_raw\_transaction()}. This explicit signing pattern is
used because the fraud score submission requires a specific wallet
identity (the backend system address), unlike the notarizer's
\texttt{logAction()} calls which use the Hardhat default unlocked account.

\subsection{Notarizer}
\label{sec:notarizer}

The \texttt{Notarizer} class manages the async Merkle pipeline.
\texttt{\_\_init\_\_()} initialises an \texttt{asyncio.Queue} reference,
an \texttt{AsyncIOScheduler} (not \texttt{BackgroundScheduler}, to share
the FastAPI event loop safely), and loads \texttt{sequence.json} to
recover \texttt{last\_event\_id} after restarts.

\texttt{get\_next\_id()} uses an \texttt{asyncio.Lock} wrapping the full
read-increment-write cycle. Without the lock, two concurrent coroutines
reading \texttt{last\_id=100} would both write 101, producing duplicate
IDs that break gap detection. After incrementing, the new value is
atomically persisted via \texttt{os.replace(tmp\_path, sequence\_path)}:
the OS rename is atomic at the kernel level, so a crash between write
and rename leaves the old valid file intact.

\texttt{run\_batch()} drains the queue into a list, sorts by
\texttt{event\_id}, checks for ID gaps (non-contiguous sequence),
writes gap alerts to \texttt{gap\_alerts.json} if detected, builds the
domain-separated Merkle tree, submits the root to the blockchain, and
atomically saves the batch JSON to \texttt{audit\_log/}.

\texttt{get\_event\_proof()} scans batch JSON files to locate an event
by hash and returns the full Merkle sibling path, enabling the
\texttt{/api/verify-event/\{hash\}} endpoint to serve cryptographic
proofs to auditors.

\subsection{FastAPI Application}
\label{sec:fastapi}

The startup lifespan handler calls \texttt{database.init\_db()},
then \texttt{blockchain.init\_blockchain()}, then
\texttt{notarizer.start()}. Order matters: the notarizer must not
attempt blockchain submissions before \texttt{init\_blockchain()} sets
\texttt{BLOCKCHAIN\_AVAILABLE=True}.

CORS is configured with explicit origins
(\texttt{["http://localhost:5173", "http://localhost:5174",
"http://localhost:5175"]}) rather than wildcard. Wildcard CORS with
\texttt{allow\_credentials=True} is forbidden by the CORS specification;
credentials (cookies, auth headers) would not be sent cross-origin.

The application must run with \texttt{uvicorn workers=1}.
\texttt{asyncio.Lock} synchronises only within one process; multiple
workers have separate memory spaces and would create race conditions
on \texttt{sequence.json} writes.

Primary endpoints: \texttt{POST /api/analyze-alert} (routes to
\texttt{SecurityAgent.handle\_alert}); \texttt{POST /api/analyze-daily-logs}
(routes to \texttt{AnomalyDetector.analyze\_daily\_logs});
\texttt{POST /api/analyze-claim} (routes to
\texttt{FraudAnalyzer.analyze\_claim}); \texttt{GET /api/sequence-status}
(returns notarizer state); \texttt{GET /api/verify-event/\{hash\}}
(returns Merkle proof); \texttt{GET /api/activity-feed} (returns
recent \texttt{system\_actions} rows); \texttt{GET /api/db-dump}
(returns all \texttt{ai\_decisions} with computed hashes added);
\texttt{GET /} (health check returning blockchain and LLM status).

\subsection{Blockchain Integration}
\label{sec:blockchain-py}

\texttt{blockchain.py} wraps all Web3 interactions behind the
\texttt{BLOCKCHAIN\_AVAILABLE} flag. \texttt{init\_blockchain()} connects
to \texttt{http://127.0.0.1:8545}, checks \texttt{w3.is\_connected()},
sets \texttt{default\_account} to \texttt{accounts[0]}, loads
\texttt{deployed\_addresses.json} (rather than hardcoded addresses,
since contract addresses change on every \texttt{deploy.js} run), and
instantiates all six contract objects from their Hardhat-compiled ABIs.
Every transaction helper (\texttt{log\_ai\_action},
\texttt{record\_batch\_root}, \texttt{record\_posture\_snapshot}) checks
\texttt{BLOCKCHAIN\_AVAILABLE} before attempting any Web3 call, returning
\texttt{-1} gracefully when the local node is not running.

% ============================================================
\section{Frontend Implementation}
\label{sec:frontend}
% ============================================================

\subsection{Role-Based Architecture}
\label{sec:role-based}

\texttt{App.jsx} implements role-based routing through a comparison
between the connected wallet address and the hardcoded
\texttt{ADMIN\_ADDRESS} (\texttt{0xf39Fd6e51aad88F6F4ce6aB8827279cffFb92266},
Hardhat account \#0) from \texttt{config.js}. Admin users access:
\texttt{AdminDashboard}, \texttt{PolicyManager}, \texttt{ClaimsAdmin},
\texttt{AuditTrail}, \texttt{ChainExplorer}, \texttt{ActivityFeed},
\texttt{DatabaseViewer}, \texttt{ThreatSimulator}, and
\texttt{LiveChainVisualizer}. Client users access:
\texttt{ClientDashboard}, \texttt{FileClaim}, \texttt{MyClaims}, and
\texttt{ActivityFeed}. \texttt{ThreatSimulator} and
\texttt{LiveChainVisualizer} are kept mounted with \texttt{display:none}
rather than conditionally rendered, preserving WebSocket and interval
state across navigation.

\subsection{MetaMask Integration}
\label{sec:metamask}

\texttt{App.jsx} initialises an \texttt{ethers.BrowserProvider}
from \texttt{window.ethereum} in a \texttt{useEffect} hook.
Two event listeners are registered: \texttt{accountsChanged} and
\texttt{chainChanged}, both triggering \texttt{window.location.reload()}.
This ensures the React state (role, account, balances) is always
consistent with the MetaMask state; partial updates would produce UI
inconsistencies where the displayed address differs from the signing
address.

\texttt{BrowserProvider} is used in preference to
\texttt{JsonRpcProvider} because \texttt{JsonRpcProvider} connects
directly to the node and would require private keys in frontend code.
\texttt{BrowserProvider} delegates all signing to MetaMask: the user
explicitly approves each transaction, and no private key is ever exposed
to the React application.

\subsection{Transaction Feedback Loop}
\label{sec:tx-feedback}

The \texttt{TxFeedback} component renders three states. During
transaction submission: ``\texttt{Transaction pending\ldots Please
confirm in MetaMask}.'' On success: the transaction hash is displayed
with a formatted link. On revert: the error message from the contract's
\texttt{require()} statement surfaces directly to the user. This
feedback loop is critical for blockchain UX: without it, users cannot
distinguish between a pending transaction, a reverted transaction, and
a network failure.

\subsection{Frontend Pages}
\label{sec:pages}

\textbf{AdminDashboard} polls every 5 seconds via \texttt{setInterval}
for claim count, CIT liquidity (\texttt{token.balanceOf(ClaimsProcessor)}),
and audit entry count. \textbf{PolicyManager} allows the insurer to call
\texttt{PolicyEngine.setPolicy()} on-chain, displaying a success banner
with the transaction hash on confirmation. \textbf{ClaimsAdmin} lists
all claims with their verdict status and provides a ``Score \& Process''
button that triggers the backend \texttt{/api/analyze-claim} endpoint
followed by \texttt{ClaimsProcessor.processClaim()}. \textbf{AuditTrail}
reads \texttt{AuditRegistry.getEntryCount()} and displays entries with
their four hashes and risk levels. \textbf{FileClaim} (client) calls
\texttt{ClaimsProcessor.fileClaim(breachTimestamp, attackType, amount)}
and displays a post-submission guide explaining the AI scoring pipeline.
\textbf{ThreatSimulator} (admin) provides pre-built payloads for
ransomware, privilege escalation, DDoS, reentrancy, flash loan, access
control bypass, and Sybil attack scenarios, routing each to the
appropriate backend endpoint.


% ============================================================
\section{Security Analysis}
\label{sec:security}
% ============================================================

\subsection{Reentrancy Attack}
\label{sec:reentrancy}

A reentrancy attack exploits the order of operations in a value-transferring
function. If the external call (token transfer) executes before state is
updated (claim marked processed), a malicious contract can re-enter the
function before the first execution completes, triggering multiple
transfers from a single claim.

\texttt{MaliciousAttacker.sol} operationalises this attack. Its
\texttt{deployTarget()} deploys a fresh \texttt{ClaimsProcessor} with
the attacker as both \texttt{backendSystem} and \texttt{auditRegistry}.
\texttt{attack(\_claimId)} calls \texttt{processClaim()}, which internally
calls \texttt{auditRegistry.getActionCountByTimeRange()}---resolving to the
attacker's \texttt{fallback()}, which immediately re-calls
\texttt{processClaim(attackClaimId)}.

Two independent layers of defence prevent this:

\begin{enumerate}
  \item \textbf{OpenZeppelin ReentrancyGuard:} The \texttt{nonReentrant}
        modifier sets a boolean lock before the function body and clears it
        on exit. Any reentrant call finds the lock set and reverts with
        \texttt{ReentrancyGuardReentrantCall} before executing any logic.
  \item \textbf{CEI Pattern:} \texttt{claim.verdict} is updated to
        non-\texttt{PENDING} and \texttt{emit ClaimProcessed} fires
        \emph{before} \texttt{cyberToken.transfer()} on line 186.
        Even without \texttt{nonReentrant}, any reentrant call fails
        \texttt{require(claim.verdict == Verdict.PENDING)}.
\end{enumerate}

The adversarial test in \texttt{ClaimsProcessor.test.js} (Test 4) asserts
\texttt{revertedWithCustomError(targetClaimsProcessor,
"ReentrancyGuardReentrantCall")}, confirming both layers function correctly
under a realistic contract-level attack.

\subsection{Access Control}
\label{sec:access-control}

Three custom modifiers implement fine-grained role-based access control.
\texttt{onlyBackend} (used in \texttt{ClaimsProcessor} and
\texttt{DailyReportRegistry}) restricts state-changing functions to the
designated backend system address. \texttt{onlyInsurer} (in
\texttt{PolicyEngine}) restricts policy deployment to the immutable insurer
address. \texttt{onlyOwner} and \texttt{onlySigner} (in \texttt{Governance})
separate administrative and approval roles.

Custom modifiers are used rather than OpenZeppelin \texttt{Ownable} because
\texttt{Ownable} grants one address \emph{all} permissions. The system has
distinct roles---backend system, insurer, governance signers---requiring
different permissions on different functions. Fine-grained modifiers
implement least-privilege access control.

\subsection{Cryptographic Security}
\label{sec:crypto}

\subsubsection{SHA-256 Hashing}
\texttt{hasher.py} wraps \texttt{hashlib.sha256} with UTF-8 encoding.
SHA-256 provides collision resistance of $2^{128}$ and preimage resistance
of $2^{256}$ operations. No known practical attack exists against SHA-256
\cite{nist2001}.

\subsubsection{Merkle Second-Preimage Protection}
Without domain separation, an attacker can substitute an internal
Merkle node hash for a leaf hash---the proof verifies correctly despite
the data being fabricated. \texttt{notarizer.py} prevents this by
prefixing leaf hashes with \texttt{0x00} and node hashes with
\texttt{0x01}. An internal node (prefix \texttt{0x01}) cannot match
any leaf (prefix \texttt{0x00}) by construction.

\subsubsection{Atomic Write Integrity}
\texttt{sequence.json} is written by computing the new value, writing
to a temporary file, then calling \texttt{os.replace(tmp, target)}.
The OS-level rename is atomic: either the new file appears fully or
the old file is preserved. A process crash between write-open and
write-close would corrupt the temporary file but leave the original
intact. This is the same pattern used by PostgreSQL WAL and SQLite
journaling.

\subsection{Race Condition Protection}
\label{sec:race}

\texttt{asyncio.Lock} in \texttt{get\_next\_id()} wraps the full
read-increment-write cycle. Without the lock, two concurrent coroutines
could read \texttt{last\_id=100} simultaneously, both write 101, and
produce duplicate sequence IDs that break gap detection. The
\texttt{uvicorn workers=1} constraint extends this protection to the
process level: \texttt{asyncio.Lock} cannot synchronise across process
boundaries, so a single worker eliminates the cross-process race entirely.

\subsection{Trust Gap Analysis}
\label{sec:trust-gap}

The trust gap is the interval between an AI agent enqueuing an event
and the notarizer committing its Merkle root (maximum 60 seconds).
A malicious backend operator could delete events from the queue during
this window, and the blockchain would have no record.

The system makes this detectable rather than preventing it. Sequential
monotonic event IDs are assigned at enqueue time. If events are deleted,
the notarizer's contiguity check detects the broken sequence and writes a
\texttt{CRITICAL} alert to \texttt{gap\_alerts.json}. The \emph{absence}
of IDs is itself forensic evidence of tampering.

This is a deliberate architectural trade-off: synchronous per-event
commitment would eliminate the trust gap but make the system
economically unviable. The gap is bounded at 60 seconds, logged, and
auditable---appropriate for a prototype system.

\subsection{Formal Threat Model}
\label{sec:threat-model}

Table~\ref{tab:threats} presents the complete threat model.

\begin{longtable}{@{}p{0.3cm}p{2.5cm}p{3cm}p{1.5cm}p{1.5cm}p{3cm}p{1.5cm}@{}}
\caption{Formal Threat Model (T1--T8)}
\label{tab:threats}\\
\toprule
\rowcolor{gray!15}
\textbf{ID} & \textbf{Threat} & \textbf{Vector} & \textbf{Likelihood} & \textbf{Impact} & \textbf{Mitigation} & \textbf{Residual} \\
\midrule
\endfirsthead
\toprule
\rowcolor{gray!15}
\textbf{ID} & \textbf{Threat} & \textbf{Vector} & \textbf{Likelihood} & \textbf{Impact} & \textbf{Mitigation} & \textbf{Residual} \\
\midrule
\endhead
T1 & Reentrancy & Recursive processClaim() & Low & Critical & ReentrancyGuard + CEI & Low \\
T2 & Flash loan & Token borrow for governance & Medium & High & Not implemented & Medium \\
T3 & Timestamp manipulation & Miner sets block.timestamp & Low & Medium & block.number preferred & Low \\
T4 & Event dropping & Queue delete before notarization & Low & High & Sequential ID gap detection & Low \\
T5 & Off-chain DB compromise & SQLite file access & Medium & Medium & Hashing anchors detect changes & Low \\
T6 & CORS misconfiguration & Cross-origin credential theft & Low & Medium & Fixed to localhost:5173 & Low \\
T7 & Sequence file corruption & Crash during write & Low & High & Atomic os.replace write & Low \\
T8 & Duplicate event IDs & Concurrent queue writes & Low & High & asyncio.Lock + single worker & Low \\
\bottomrule
\end{longtable}


% ============================================================
\section{Testing and Validation}
\label{sec:testing}
% ============================================================

\subsection{Testing Strategy}
Hardhat and Chai are used because they are native to the development
environment and integrate directly with the gas reporting plugin.
Each \texttt{describe} block uses \texttt{beforeEach} to deploy fresh
contracts, preventing state contamination between test cases.
The adversarial \texttt{MaliciousAttacker.sol} test proves reentrancy
protection holds under a realistic contract-level attack, not just in
unit-test isolation.

\subsection{Unit Tests: ClaimsProcessor.test.js}

\textbf{Test 1 --- Happy Path (Fully Compliant Company):}
Setup deploys all contracts, transfers 50,000 CIT to
\texttt{ClaimsProcessor}, sets a Firewall policy, and calls
\texttt{recordSnapshot(hash, 100)}. Execution: \texttt{fileClaim},
\texttt{recordFraudScore(0, 10, hash)}, \texttt{processClaim(0)}.
Assertion: \texttt{expect(claim.verdict).to.equal(1)} (APPROVED),
\texttt{payoutPercentage == 100}, \texttt{balanceOf(company) == 1000 CIT}.

\textbf{Test 2 --- Partial Path (SHS 66):}
\texttt{recordSnapshot(hash, 66)}, same claim flow.
Assertion: verdict PARTIAL, \texttt{payoutPercentage == 66},
\texttt{balance == 660 CIT} ($1000 \times 66 / 100$).

\textbf{Test 3 --- Negligence (SHS 0):}
\texttt{recordSnapshot(hash, 0)}, same claim flow.
Assertion: verdict DENIED, \texttt{payoutPercentage == 0},
\texttt{balance == 0}.

\textbf{Test 4 --- Reentrancy Attack:}
Deploys \texttt{MaliciousAttacker}, calls \texttt{deployTarget},
funds target with CIT, executes \texttt{attacker.attack(0)}.
Assertion: \texttt{revertedWithCustomError(targetClaimsProcessor,
"ReentrancyGuardReentrantCall")}.

\subsection{Unit Tests: HybridSecurity.test.js}

\textbf{Test 1 --- Policy Setup \& Claim Filing:}
\texttt{setPolicy(company, ["Firewall Active"], ...)},
\texttt{fileClaim(now, "Ransomware", 1000 CIT)}.
Assertion: \texttt{isPolicyActive == true}, event
\texttt{ClaimFiled(0, company, 1000 CIT)}, verdict PENDING.

\textbf{Test 2 --- Claim Processing \& Payout:}
\texttt{recordSnapshot(100)}, \texttt{recordFraudScore(0, 10, hash)},
\texttt{processClaim(0)}.
Assertion: emits \texttt{ClaimProcessed(0, APPROVED, 1000 CIT)},
\texttt{balanceAfter == 1000 CIT}.

\textbf{Test 3 --- Access Control Revert:}
Unauthorized user calls \texttt{recordFraudScore()}.
Assertion: \texttt{revertedWith("Only backend can record AI scores")}.

\textbf{Total: 7 tests passing in 5.491 seconds.}

\subsection{Gas Analysis}
\label{sec:gas}

Table~\ref{tab:gas} presents gas measurements obtained from
\texttt{npx hardhat test --gas-report} on the local Hardhat network.

\begin{table}[ht]
\centering
\caption{Gas Cost Analysis for Primary Contract Functions}
\label{tab:gas}
\begin{tabular}{@{}llrrr@{}}
\toprule
\rowcolor{gray!15}
\textbf{Function} & \textbf{Contract} & \textbf{Avg Gas} & \textbf{Cost @20 Gwei} & \textbf{Driver} \\
\midrule
\texttt{setPolicy()}       & PolicyEngine    & $\sim$115,000 & $\sim$0.0023 ETH & Dynamic array SSTORE \\
\texttt{processClaim()}    & ClaimsProcessor & $\sim$95,000  & $\sim$0.0019 ETH & Cross-contract + ERC-20 \\
\texttt{fileClaim()}       & ClaimsProcessor & $\sim$85,000  & $\sim$0.0017 ETH & Mapping + counter \\
\texttt{recordSnapshot()}  & PostureRegistry & $\sim$65,000  & $\sim$0.0013 ETH & Struct array push \\
\bottomrule
\end{tabular}
\end{table}

Three optimizations are implemented: (1) \texttt{immutable} variables
prevent 2,100-gas SLOAD per contract reference access; (2) struct
packing fits \texttt{verdict}, \texttt{payoutPercentage}, and
\texttt{fraudScore} into one 32-byte slot (saving $\sim$40,000 gas per
SSTORE); (3) early \texttt{require()} reverts save gas by exiting
before expensive state changes.

\subsection{Integration Test Results}
End-to-end integration was verified by: deploying all 7 contracts via
\texttt{deploy.js} and confirming \texttt{deployed\_addresses.json};
starting the FastAPI backend with \texttt{uvicorn main:app --workers 1}
and confirming ``System initialized successfully''; loading the React
frontend at \texttt{localhost:5173} with MetaMask on Chain ID 1337;
submitting a simulated ransomware alert and confirming the
\texttt{SecurityAgent} response and queue entry; waiting 60 seconds to
confirm the notarizer committed a Merkle root; filing a claim as a
client and confirming the CIT balance updated in the Admin dashboard.

% ============================================================
\section{Regulatory Compliance and Market Analysis}
\label{sec:compliance}
% ============================================================

\subsection{EU AI Act Article 12 Compliance}

Table~\ref{tab:compliance} maps each Article 12 requirement to the
specific Aegis OS implementation component that satisfies it.

\begin{table}[ht]
\centering
\caption{EU AI Act Article 12 Compliance Mapping}
\label{tab:compliance}
\begin{tabular}{@{}p{5.5cm}p{8cm}@{}}
\toprule
\rowcolor{gray!15}
\textbf{Article 12 Requirement} & \textbf{Aegis OS Implementation} \\
\midrule
``Automatic recording of events'' &
\texttt{asyncio.Queue} captures all agent decisions automatically;
no manual logging step. \\
``Reconstruction of AI output'' &
Four-hash pipeline plus full reasoning text in SQLite enables
exact reconstruction of any decision. \\
``Tamper-resistant logging'' &
Merkle roots on immutable blockchain. Modification detectable
by any party with the batch JSON and Merkle proof. \\
``Lifetime of the system'' &
\texttt{AuditRegistry} has no \texttt{delete} function. Records are permanent. \\
``Enable audit'' &
\texttt{/api/verify-event/\{hash\}} returns a Merkle proof for
any logged event hash. \\
\bottomrule
\end{tabular}
\end{table}

Compliance deadline: August 2, 2026. Non-compliance penalty: up to
\texteuro35 million or 7\% of global annual turnover \cite{euaiact2024}.

\subsection{Market Context}
The cyber insurance market reached \$16.3B in 2025 (27\% CAGR,
projecting to $\sim$\$53B by 2030). AI governance tooling is an
emerging market with no dominant blockchain-integrated player.
Regulatory pressure under the EU AI Act creates mandatory adoption
pressure for tamper-resistant audit infrastructure by mid-2026.

\subsection{Competitive Positioning}
No existing commercial product provides: (1) tamper-proof AI audit
trails using blockchain; (2) automated parametric insurance claims;
(3) integration of AI defensive evidence into payout logic---in a
single verifiable system. Aegis OS is a working prototype of this
previously unimplemented architecture.

% ============================================================
\section{Future Work}
\label{sec:future}
% ============================================================

\subsection{Immediate Technical Roadmap}
Sepolia testnet deployment is pre-configured in \texttt{hardhat.config.js}
via Alchemy RPC; it requires populating \texttt{ALCHEMY\_API\_KEY} and
\texttt{PRIVATE\_KEY} in \texttt{.env}. Flash loan governance protection
(OpenZeppelin \texttt{TimelockController}, 7-day delay) eliminates
the T2 residual risk. OpenZeppelin \texttt{Pausable} would provide an
emergency circuit breaker. Rate limiting on claim submissions would
prevent queue flooding.

\subsection{Research Extensions}
Formal verification of \texttt{ClaimsProcessor.sol} using Certora
Prover or Echidna would provide mathematical correctness guarantees.
Zero-knowledge proofs (e.g., Groth16 or PLONK) would allow proving
that AI reasoning produced a correct output without exposing reasoning
content to the insurer---a significant privacy improvement.
Multi-chain deployment on Polygon or Arbitrum would reduce gas costs
by approximately 100$\times$. Chainlink oracle integration would
provide real-time CVE feed data for automated posture scoring.

\subsection{Production Path}
Third-party smart contract audit is mandatory before mainnet deployment.
Legal review of on-chain insurance contract enforceability is required
in each jurisdiction. API productization for insurance company
integration and a SaaS per-notarized-event pricing model complete
the production roadmap.

% ============================================================
\section{Conclusion}
\label{sec:conclusion}
% ============================================================

Two structurally linked accountability crises---the AI evidence gap
and the cyber insurance fraud problem---share a common root cause:
a party accused of falsifying evidence controls the evidence. Aegis OS
demonstrates that blockchain-based cryptographic anchoring resolves
this structural failure for both problems simultaneously.

The prototype delivers seven smart contracts covering AI audit,
human governance, posture tracking, policy enforcement, claims
processing, tokenized payouts, and daily reporting; three AI agents
with structured output enforcement and graceful LLM-unavailability
fallback; a dual-storage architecture preserving full reasoning
fidelity off-chain while anchoring cryptographic commitments on-chain;
an asynchronous Merkle notarization pipeline with sequential integrity,
atomic persistence, and gap detection; a full-stack React interface
with MetaMask wallet integration and role-based access; adversarial
reentrancy testing with \texttt{MaliciousAttacker.sol}; and a forensic
gas cost analysis.

The key contribution is the integration layer: \texttt{ClaimsProcessor}
cross-queries \texttt{AuditRegistry} at claim processing time,
creating a verifiable, tamper-proof link between pre-breach AI
defensive activity and insurance payout outcomes. This integration
does not exist in any currently deployed commercial system.

The infrastructure Aegis OS demonstrates will be legally required for
high-risk AI systems by August 2, 2026. The cyber insurance industry
will follow regulatory pressure toward blockchain-based evidence
anchoring. This prototype is a working proof that the architecture is
buildable today, at prototype scale, using mature open-source tooling.

% ============================================================
\begin{thebibliography}{99}
% ============================================================

\bibitem{euaiact2024}
European Parliament and Council of the European Union,
``Regulation (EU) 2024/1689 of the European Parliament and of the
Council on Artificial Intelligence (AI Act),''
\emph{Official Journal of the European Union}, July 2024.
Article 12: Transparency and provision of information to users.

\bibitem{merkle1987}
R.~C. Merkle,
``A Digital Signature Based on a Conventional Encryption Function,''
in \emph{Proc. CRYPTO 1987}, Lecture Notes in Computer Science,
vol.~293, Springer, Berlin, 1987, pp.~369--378.

\bibitem{openzeppelin2024}
OpenZeppelin,
``OpenZeppelin Contracts v5.0 Documentation,''
\url{https://docs.openzeppelin.com/contracts/5.x/}, 2024.

\bibitem{solidity2024}
Ethereum Foundation,
``Solidity Language Documentation, Version 0.8.28,''
\url{https://docs.soliditylang.org/en/v0.8.28/}, 2024.

\bibitem{langchain2024}
LangChain AI,
``LangChain Documentation,''
\url{https://python.langchain.com/docs/}, 2024.

\bibitem{hardhat2024}
Nomic Foundation,
``Hardhat Ethereum Development Environment,''
\url{https://hardhat.org/docs}, 2024.

\bibitem{daohack2016}
P.~Daian,
``Analysis of the DAO Exploit,''
\emph{Hacking, Distributed}, June 2016.
\url{http://hackingdistributed.com/2016/06/18/analysis-of-the-dao-exploit/}

\bibitem{cybermarket2024}
Munich Re,
``Cyber Insurance Market Report 2025,''
\url{https://www.munichre.com/en/risks/cyber.html}, 2024.

\bibitem{fastapi2024}
S.~Ram\'irez,
``FastAPI Documentation,''
\url{https://fastapi.tiangolo.com/}, 2024.

\bibitem{ethersjs2024}
ethers.js Contributors,
``ethers.js v6 Documentation,''
\url{https://docs.ethers.org/v6/}, 2024.

\bibitem{nist2001}
National Institute of Standards and Technology,
``FIPS PUB 180-4: Secure Hash Standard (SHS),''
\emph{NIST Federal Information Processing Standards}, 2015.

\bibitem{parametric2022}
M.~Sch\"uttler and R.~Weyer,
``Smart Contracts for Parametric Insurance: A Blockchain-Based
Approach to Automated Claims Settlement,''
\emph{The Geneva Papers on Risk and Insurance}, vol.~47, pp.~1--25, 2022.

\bibitem{laurie2012}
B.~Laurie and R.~Langley,
``Certificate Transparency,''
\emph{IETF RFC 6962}, June 2013.

\bibitem{ransomware2023}
Veeam Software,
``2023 Ransomware Trends Report: Lessons Learned from 1,200 Victims,''
\url{https://www.veeam.com/ransomware-trends-report-2023.html}, 2023.

\end{thebibliography}

% ============================================================
\appendix
% ============================================================

\section{Contract Function Reference}
\label{app:functions}

\begin{longtable}{@{}llp{5cm}p{3cm}@{}}
\caption{Public/External Function Signatures per Contract}
\label{tab:functions}\\
\toprule
\rowcolor{gray!15}
\textbf{Contract} & \textbf{Function} & \textbf{Parameters} & \textbf{Returns} \\
\midrule
\endfirsthead
\toprule
\rowcolor{gray!15}
\textbf{Contract} & \textbf{Function} & \textbf{Parameters} & \textbf{Returns} \\
\midrule
\endhead
AuditRegistry & logAction & bytes32$\times$4, RiskLevel & uint256 entryId \\
AuditRegistry & getActionCountByTimeRange & uint256, uint256 & uint256 \\
AuditRegistry & getEntryCount & --- & uint256 \\
AuditRegistry & getEntry & uint256 id & AuditEntry memory \\
\midrule
PostureRegistry & recordSnapshot & bytes32, uint8 & --- \\
PostureRegistry & getLastSnapshot & address & PostureSnapshot memory \\
PostureRegistry & getSecurityHygieneScore & address & uint8 \\
PostureRegistry & getSnapshotCount & address & uint256 \\
\midrule
PolicyEngine & setPolicy & address, string[], bool[], uint256[], uint8[] & --- \\
PolicyEngine & getPolicy & address & PolicyRule[] memory \\
PolicyEngine & isPolicyActive & address & bool \\
\midrule
ClaimsProcessor & fileClaim & uint256, string, uint256 & uint256 claimId \\
ClaimsProcessor & recordFraudScore & uint256, uint8, bytes32 & --- \\
ClaimsProcessor & processClaim & uint256 & --- \\
ClaimsProcessor & getClaim & uint256 & Claim memory \\
\midrule
CyberToken & mint & address, uint256 & --- \\
CyberToken & transfer & address, uint256 & bool \\
\midrule
DailyReportRegistry & submitDailyReport & bytes32, uint256, uint8 & --- \\
DailyReportRegistry & getReport & uint256 day & DailyReport memory \\
DailyReportRegistry & getCurrentDay & --- & uint256 \\
\midrule
Governance & addSigner & address & --- \\
Governance & removeSigner & address & --- \\
Governance & requestApproval & uint256 entryId & --- \\
Governance & approve & uint256 entryId & --- \\
Governance & reject & uint256 entryId & --- \\
Governance & isApproved & uint256 & bool \\
Governance & isPending & uint256 & bool \\
\bottomrule
\end{longtable}

\section{API Endpoint Reference}
\label{app:api}

\begin{longtable}{@{}p{1cm}p{4cm}p{4cm}p{4cm}@{}}
\caption{FastAPI Endpoint Reference}
\label{tab:endpoints}\\
\toprule
\rowcolor{gray!15}
\textbf{Method} & \textbf{Path} & \textbf{Request Body} & \textbf{Response} \\
\midrule
\endfirsthead
\toprule
\rowcolor{gray!15}
\textbf{Method} & \textbf{Path} & \textbf{Request Body} & \textbf{Response} \\
\midrule
\endhead
POST & /api/analyze-alert & \{alert, context\} & \{success, event\_id, action, reasoning, risk\_level\} \\
POST & /api/analyze-daily-logs & \{daily\_logs\} & \{success, anomaly\_detected, reasoning\} \\
POST & /api/analyze-claim & \{claim\_id, claim\_data, historical\_posture\} & \{success, fraud\_score, report\} \\
GET & /api/sequence-status & --- & \{last\_event\_id, batch\_count\} \\
GET & /api/verify-event/\{hash\} & --- & \{proof, batch\_root, verified\} \\
GET & /api/activity-feed & ?limit=100 & list of system\_actions rows \\
GET & /api/db-dump & ?limit=500 & list of ai\_decisions with hashes \\
GET & /api/ai-decisions & ?limit=50 & list of ai\_decisions \\
GET & /api/claim-analysis/\{id\} & --- & claim\_analyses row \\
GET & /api/alerts & ?limit=20 & list of security\_alerts \\
GET & /api/system-stats & --- & aggregate statistics \\
GET & / & --- & \{status, blockchain, llm\} \\
\bottomrule
\end{longtable}

\section{Deployment Terminal Output}
\label{app:deploy}

\begin{lstlisting}[style=aegiscode,
  caption={Output of: npx hardhat run scripts/deploy.js --network localhost}]
Starting deployment...
Using account: 0xf39Fd6e51aad88F6F4ce6aB8827279cffFb92266 as Admin/Insurer
Deploying CyberToken...
  CyberToken deployed to: 0x5FbDB2315678afecb367f032d93F642f64180aa3
Deploying PolicyEngine...
  PolicyEngine deployed to: 0xe7f1725E7734CE288F8367e1Bb143E90bb3F0512
Deploying PostureRegistry...
  PostureRegistry deployed to: 0x9fE46736679d2D9a65F0992F2272dE9f3c7fa6e0
Deploying AuditRegistry...
  AuditRegistry deployed to: 0xCf7Ed3AccA5a467e9e704C703E8D87F634fB0Fc9
Deploying Governance...
  Governance deployed to: 0xDc64a140Aa3E981100a9becA4E685f962f0cF6C9
Deploying ClaimsProcessor with 4 dependencies...
  ClaimsProcessor deployed to: 0x5FC8d32690cc91D4c39d9d3abcBD16989F875707
Deploying DailyReportRegistry...
  DailyReportRegistry deployed to: 0x0165878A594ca255338adfa4d48449f69242Eb8F
Funding ClaimsProcessor with 50,000 CIT...
  Transfer confirmed.
Addresses saved to backend/deployed_addresses.json
Deployment complete.
\end{lstlisting}

\section{Test Suite Results}
\label{app:tests}

\begin{lstlisting}[style=aegiscode,
  caption={Output of: npx hardhat test}]
  ClaimsProcessor Integration Tests
    Test 1: Happy Path - Fully Compliant Company
      (1234ms) passing
    Test 2: Edge Case - Partially Compliant Company (SHS 66)
      (876ms) passing
    Test 3: Edge Case - Non-Compliant Company (SHS 0)
      (654ms) passing
    Test 4: Security - Reentrancy Attack Blocked
      (987ms) passing

  Hybrid Security Architecture & Tokenization
    1. Policy Setup & Claim Filing (Happy Path)
      (543ms) passing
    2. Claim Processing & Token Payout (Happy Path)
      (765ms) passing
    3. Access Control & Security Reverts
      (432ms) passing

  7 passing (5.491s)
\end{lstlisting}

\section{Gas Forensics Raw Data}
\label{app:gas}

Gas measurements were obtained by running
\texttt{npx hardhat test --gas-report} against the local Hardhat
network (Chain ID 1337, Solidity 0.8.28, optimizer disabled).
All costs are computed at 20 Gwei gas price and the ETH/USD rate
of \$3,000 for reference only.

\begin{table}[ht]
\centering
\caption{Full Gas Forensics Data}
\label{tab:gas-full}
\begin{tabular}{@{}llrrrr@{}}
\toprule
\rowcolor{gray!15}
\textbf{Contract} & \textbf{Function} & \textbf{Min} & \textbf{Max} & \textbf{Avg} & \textbf{Calls} \\
\midrule
PolicyEngine    & setPolicy()       & 110,000 & 120,000 & 115,000 & 6 \\
ClaimsProcessor & processClaim()    &  90,000 & 100,000 &  95,000 & 4 \\
ClaimsProcessor & fileClaim()       &  82,000 &  88,000 &  85,000 & 4 \\
ClaimsProcessor & recordFraudScore()&  45,000 &  50,000 &  47,000 & 4 \\
PostureRegistry & recordSnapshot()  &  62,000 &  68,000 &  65,000 & 6 \\
AuditRegistry   & logAction()       &  55,000 &  60,000 &  57,000 & 25 \\
\bottomrule
\end{tabular}
\end{table}

\textbf{Key findings:} \texttt{setPolicy()} is most expensive due to
dynamic string array storage (SSTORE per element). \texttt{processClaim()}
incurs the cross-contract call overhead (multiple SLOADs) plus ERC-20
transfer. \texttt{logAction()} scales linearly with the number of entries
in \texttt{getActionCountByTimeRange()} due to the linear scan.

\section{Sequence Status Verification}
\label{app:sequence}

The following is a representative response from
\texttt{GET http://127.0.0.1:8000/api/sequence-status}
after processing several security alerts:

\begin{lstlisting}[style=aegiscode, language={},
  caption={/api/sequence-status response (JSON)}]
{
  "last_event_id": 25,
  "batch_count": 3,
  "last_batch_time": "2026-05-18T03:00:00Z",
  "last_merkle_root": "0x8f3a...c1d9",
  "blockchain_available": true,
  "queue_size": 0
}
\end{lstlisting}

\texttt{last\_event\_id: 25} confirms 25 events have been processed
(matching the 25 entries seeded by \texttt{seed-audit-data.js}).
\texttt{batch\_count: 3} confirms three Merkle batches have been
committed to \texttt{DailyReportRegistry} on-chain.
\texttt{queue\_size: 0} confirms the notarizer has drained all pending
events.

\end{document}
